\section{Tail four and the quadratic growth of the triple quotient genus}
\label{sec:factorwise-four-tail}

\paragraph{Status and exact scope.}
These are root statements HD-FT09--10, reserved question Q2dh.
The complete proof and source application have passed root and separate
independent whole-proof review. The initial text, two precision edits
and the complete independent audit are retained. No finite certificate
was needed, executed or claimed for these universal inequalities.
The local character estimate is effective. The final exclusion is qualitative,
because the cited genus-growth constants have not been evaluated.
It is not an effective detector-degree bound and is not substituted into FQ.

Use the actual quotient and tensor factorization (FT.1)--(FT.2): a connected
finite-group cover of the sphere, an actual quotient onto $P=A_n$ or $S_n$,
and an actual factor $V=W\otimes U$ up to scalars, where the genuine inertia
operators satisfy
\[
 e=\dim V,\qquad \sum_i\operatorname{rank}(1-T_i)=2e+6.
 \tag{FT4.1}
\]
The partner $U$, its degree and its field are arbitrary. Let $\mu\vdash4$,
and let $W=S^{(n-4,\mu)}$ or its irreducible ordinary restriction to $A_n$.
Transposition of the partition is permitted: its sign twist does not affect
absolute character ratios or the maximum-eigenspace support.
As in D-AFT01, write
\[
 q_\mu=CX_1^4+R,\qquad C=f^\mu/24,\qquad
 L=\|R\|_1,\qquad M=1+L/C\ge1.
 \tag{FT4.2}
\]
The exact character polynomial has weighted degree at most four, assigning
weight $j$ to $X_j$, and $R$ has ordinary degree at most three.
These are the accepted consequences of Garsia--Goupil,
Proposition I.1 and its complete proof.\cite{HD:GarsiaGoupil2009}

\subsection*{HD-FT09. A localized quadratic-order error}

\begin{lemma}[HD-FT09]
If $n\ge256M$, then $a=\dim W\ge2n^3$ and, writing $f=X_1(x)$,
\[
 \frac{|\chi_W(x)|}{a}\le
 \begin{cases}
 (f/n)^3,&f\ge n/2,\\
 (f/n)^3+4M^2/n^2,&f<n/2.
 \end{cases}
 \tag{FT4.3}
\]
The estimate holds for every permutation, hence for every power of every
projected inertia.
\end{lemma}

\begin{proof}
Put $t=f/n$ and $m=n-f$. The ordinary-degree bound gives
\[
 a\ge Cn^4(1-M/n)\ge n^4/48\ge2n^3,
 \tag{FT4.4}
\]
because $C\ge1/24$ and $n\ge256>96$. In particular $n\ge256>9$, so the
stable partition is non-self-conjugate and all character premises hold.

If $f\ge n/2$, the character is positive: the leading term is at least
$Cn^4/16$, whereas the absolute remainder is at most $Ln^3$, and
$n\ge256M>16L/C$.
For a pure monomial in $f$, the difference of its value at $f$ and $n$
is at most its degree times $mn^{\deg-1}$. Every other monomial has a
non-$f$ cycle count at most $m$. Thus
\[
 |R(f,X_2,\ldots)-R(n,0,\ldots)|\le4Lm n^2.
 \tag{FT4.5}
\]
Since $1-t^3\le3m/n$ and $t^3-t^4\ge(1-t)/8$ in this range,
\[
 q_\mu-a t^3\le-Cmn^3/8+7Lm n^2\le0,
 \tag{FT4.6}
\]
using $n\ge256M>56L/C$. This includes $m=0$, when the difference is
exactly zero, and proves the first case with its absolute value.

Suppose $0\le t<1/2$. For a remainder monomial, let $j$ be the exponent
of $f$ and $b$ the sum of all other exponents. Then $j+2b\le4$ and
$j+b\le3$. After division by $n^4$, the monomial is at most $t^3/n$
if $j\ge3$, at most $t^2/n$ if $j=2$, at most $t/n^2$ if $j=1$,
and at most $n^{-2}$ if $j=0$. Coefficientwise summation therefore gives
\[
 \frac{|R|}{Cn^4}
 \le M\left(\frac{t^3}{n}+\frac{t^2}{n}
                       +\frac{t}{n^2}+\frac1{n^2}\right).
 \tag{FT4.7}
\]
The elementary arithmetic--geometric mean inequalities
$3Az^2\le2z^3+A^3$ and $3Az\le z^3+2A^{3/2}$, applied with
$A=4M$ and $z=nt$, give
\[
 \frac{Mt^2}{n}\le\frac{t^3}{6}+\frac{16M^3}{3n^3},
 \qquad
 \frac{Mt}{n^2}\le\frac{t^3}{12}+\frac{4M^2}{3n^3}.
 \tag{FT4.8}
\]
In the second inequality we used $M^{3/2}\le M^2$.
Since $t^4\le t^3/2$, equations (FT4.7)--(FT4.8) yield
\[
 \frac{|q_\mu|}{Cn^4}
 \le\left(\frac34+\frac Mn\right)t^3
       +\frac{M}{n^2}+\frac{16M^3+4M^2}{3n^3}.
 \tag{FT4.9}
\]
Divide by $a/(Cn^4)\ge1-M/n$.
The cubic coefficient is at most one, since $M/n\le1/256<1/8$.
Moreover
\[
 M+\frac{16M^3+4M^2}{3n}
 \le M^2+\frac{20M^3}{3\cdot256M}\le2M^2,
 \qquad (1-M/n)^{-1}\le2.
 \tag{FT4.10}
\]
The remaining term is consequently at most $4M^2/n^2$.
This proves the second case of (FT4.3).
\end{proof}

\subsection*{HD-FT10. A qualitative cutoff from a published genus-growth estimate}

We state precisely the additional published input. In the proof of
Monderer--Neftin Theorem 2.5, equation (2.2), there are positive constants
$c,d$ such that, for the genera $g_j$ of the quotients by unordered
$j$-set stabilizers of an actual $A_n$ or $S_n$ sphere cover,
\[
 g_j-g_{j-1}>(cn-dj^{15})
               \frac{\binom nj}{\binom n2}\qquad(3\le j\le\lfloor n/2\rfloor).
 \tag{FT4.11}
\]
The $j\ge3$ range applies also to the exceptional ramification types
in their Table 2.1.\cite[proof of Theorem~2.5, equation~(2.2)]{HD:MondererNeftin2022}
We may include an initial threshold $n_0\ge20$ in this asymptotic
form without strengthening the source. No numerical $c,d,n_0$ is claimed.
This is the estimate expressly invoked in the published proof, read in
the author's revised 18 December 2020 text. We do not claim a fresh
proof of its cited Neftin--Zieve precursor, nor an independently computed
effective version of those constants.

For the fixed $\mu$, define the integer, depending on these unevaluated
constants,
\[
 N_\mu^{(4)}=
 \left\lceil\max\left\{n_0,256M,4,
                  \frac{2d\,3^{15}}c,\frac{80M^2}c\right\}\right\rceil.
 \tag{FT4.12}
\]

\begin{theorem}[HD-FT10]
No source in the actual factorwise setup (FT4.1) with the indicated
tail-four factor exists for $n>N_\mu^{(4)}$.
Thus all five tail-four labels have a finite qualitative natural-degree
cutoff, with arbitrary tensor partners, scalars and coefficient fields.
This statement does not specify a numerical cutoff or an effective
enumeration range.
\end{theorem}

\begin{proof}
Put $\varepsilon=4M^2/n^2$. Call an inertia image bad if it fixes fewer
than $n/2$ letters. A nonbad image has only nonbad powers. Fourier inversion,
the cube fixed-point formula and tensor spectral convolution give
\[
 \frac{\operatorname{rank}(1-T_i)}e
 \ge\frac{\operatorname{ind}(x_i,\{1,\ldots,n\}^3)}{n^3}
                -\varepsilon\mathbf1_{\{i\text{ bad}\}}.
 \tag{FT4.13}
\]
Identity quotient images cause no problem: their index and error are zero.
There are at most ten bad entries. Indeed the same Fourier bound, with
$u^3\le u$ for $0\le u\le1$, gives at a bad entry
\[
 \frac{\nu(W(x_i))}{a}
 \ge\frac{\operatorname{ind}(x_i)}n-\varepsilon
 >\frac14-\varepsilon>\frac15,
 \qquad \varepsilon\le\frac1{16384}<\frac1{20}.
 \tag{FT4.14}
\]
Here $\operatorname{ind}(x_i)\ge(n-f(x_i))/2>n/4$.
Eleven such entries would exceed $11e/5$, whereas (FT4.4) makes the
actual total budget $2+6/e\le2+3/n^3<11/5$.

Let $\widehat\gamma_3$ be the genus of the unordered-triple quotient.
Apply (FT4.11) with $j=3$ and use $g_2\ge0$. The threshold gives
\[
 \widehat\gamma_3>
   (cn-d\,3^{15})\frac{n-2}{3}
       \ge\frac{cn}{2}\frac n6=\frac{cn^2}{12}.
 \tag{FT4.15}
\]
The ordered-distinct-triple quotient has degree six over this curve, so
its genus satisfies $\gamma_3\ge6\widehat\gamma_3-5>cn^2/2-5$.
There are five $P$-orbits on the cube, with genera
$\gamma_1,\gamma_2,\gamma_2,\gamma_2,\gamma_3$.
The actual quotient-cover Riemann--Hurwitz identities therefore imply
\[
 \sum_i\operatorname{ind}(x_i,\{1,\ldots,n\}^3)
 =2(n^3-5)+2(\gamma_1+3\gamma_2+\gamma_3)
 >2n^3+cn^2-20.
 \tag{FT4.16}
\]
Sum (FT4.13) and subtract at most ten errors. Then
\[
 2+\frac6e>2+\frac cn-\frac{40M^2}{n^2}-\frac{20}{n^3}.
 \tag{FT4.17}
\]
But $cn\ge80M^2$ and $n\ge4$ show
\[
 cn^2-40M^2n-23\ge40M^2n-23>0.
 \tag{FT4.18}
\]
Thus the right side of (FT4.17) is strictly larger than
$2+3/n^3\ge2+6/e$, a contradiction.
\end{proof}

\paragraph{Gate and geometric interpretation.}
Together with HD-FT01--08, this leaves only ordinary tails of size at most
three as potentially unbounded in the stated actual factorwise setup.
The genus-growth input gives a qualitative reduction, not the effective
integral detector bound required in the main problem. Obtaining numerical
constants, deciding the remaining smaller tails, and supplying actual
field-compatible cycles are distinct outstanding tasks.
The argument concerns actual finite covers, their weight-one local systems
and rational deck projectors. It narrows supplier routes for integral
uniform condition A (cote16, scan050; SIGN\_NEXT SN-C03/C04), while the
condition I positivity step (scan049; cote44, Theorem4), the Cayley--Hamilton
inverse of scans048--049, and the accepted EGA/SGA/FGA spreading interfaces
do not produce missing cycles. No Markman, Severn or032 premise is added.
